\subsection{Using Composition Series}

Let A be a ring. Let E be an A-module of finite length and S a simple A-module. By the Jordan--Hölder theorem (I, §4, No.~7, p.~43, Theorem 6), the number of quotients of a Jordan--Hölder series of E that are isomorphic to S does not depend on the sequence. We denote it by \(\ell_{\mathrm{S}}(\mathrm{E})\) and call it the multiplicity of S in E. The support of the A-module E is the set of classes of simple A-modules S such that \(\ell_{\mathrm{S}}(\mathrm{E}) \neq 0\) . When E is semisimple and has finite length, the integer \(\ell_{\mathrm{S}}(\mathrm{E})\) is the length {[}E : S{]} of the isotypical component of E of type S (VIII, p.~72), and the notion of support coincides with that introduced in VIII, p.~66.

Lemma 1. --- Let E, E', and E'\,' be A-modules of finite length and

\[
0 \longrightarrow \mathrm{E} ^ {\prime} \stackrel {i} {\longrightarrow} \mathrm{E} \stackrel {p} {\longrightarrow} \mathrm{E} ^ {\prime \prime} \longrightarrow 0
\]

an exact sequence. We have \(\ell_{\mathrm{S}}(\mathrm{E}) = \ell_{\mathrm{S}}(\mathrm{E}') + \ell_{\mathrm{S}}(\mathrm{E}'')\) .

Let \(\Sigma'\) and \(\Sigma''\) be Jordan-Hölder series of \(i(\mathrm{E}')\) and \(\mathrm{E}/i(\mathrm{E}')\) , respectively. There exists a Jordan-Hölder series \(\Sigma\) of \(\mathrm{E}\) whose sequence of quotients can be obtained by juxtaposing the sequence of quotients of \(\Sigma\) and that of \(\Sigma'\) (I, §4, No.~8, p.~44).

PROPOSITION 7. --- Let C be a hereditary set of classes of modules such that every module of type C has finite length. Let S be the set of classes of simple modules belonging to C. Then the family \(([S]_{\mathcal{C}})_{S \in \mathcal{S}}\) is a basis of the Z-module \(K(\mathcal{C})\) , and we have

\[
[ \mathrm{E} ] _ {\mathcal {C}} = \sum_ {\mathrm{S} \in \mathcal {S}} \ell_ {\mathrm{S}} (\mathrm{E}) [ \mathrm{S} ] _ {\mathcal {C}}\tag{8}
\]

for every module E of type \(\mathcal{C}\) .

Formula (8) follows from Proposition 3 applied to a Jordan--Hölder series of E. By Lemma 1, for every element S of S, there exists a Z-linear mapping \(\varphi_{\mathrm{S}}\) from \(\mathrm{K}(\mathcal{C})\) to Z such that \(\varphi_{\mathrm{S}}([\mathrm{E}]_{\mathcal{C})} = \ell_{\mathrm{S}}(\mathrm{E})\) for every module E of type C. In particular, we have \(\varphi_{\mathrm{S}}([\mathrm{S}]_{\mathcal{C}}) = 1\) and \(\varphi_{\mathrm{S}}([\mathrm{S}']_{\mathcal{C}}) = 0\) for every \(S' \neq S\) in S. It follows that the elements of the form \([S]_{\mathcal{C}}\) (for \(S \in S\) ) are linearly independent over Z; these elements generate \(\mathrm{K}(\mathcal{C})\) by formula (8).

COROLLARY. --- Let E and F be semisimple modules of type C. The module E is isomorphic to F if and only if we have \([E]_{C} = [F]_{C}\) in \(K(\mathcal{C})\) .

Indeed, we have \([E]_{\mathcal{C}} = \sum_{S \in \mathcal{S}} \ell_{S}(E)[S]_{\mathcal{C}}\) and an analogous formula for F, and E is isomorphic to F if and only if we have \(\ell_{S}(E) = \ell_{S}(F)\) for every \(S \in \mathcal{S}\) (VIII, p.~72).

Remark. --- The set \(\mathrm{K}(\mathcal{C})^{+}\) is the submonoid of \(\mathrm{K}(\mathcal{C})\) generated by the family \(([S]_{\mathcal{C}})_{S\in\mathcal{S}}\) .

